\section{Computational Details and Methodological Protocol}

\subsection{Quantum MOF (QMOF) Ingestion and Electrocatalytic Screening}
The structural dataset was retrieved from the Materials Project MOF Explorer / QMOF database (\cite{qmof2021}), encompassing $>20,000$ DFT-relaxed metal--organic frameworks. To focus on electrocatalytically tractable and experimentally realizable systems, hierarchical filtering was applied:
\begin{enumerate}
    \item \textbf{Experimental Synthesizability:} Only frameworks with validated experimental synthesis records were retained ($\mathrm{synthesized} = \mathrm{True}$).
    \item \textbf{Active Transition Metal Nodes:} Filtered for redox-active $3d$, $4d$, and $5d$ transition metals ($\mathrm{Co, Cu, Fe, Mn, Mo, Ni, Ru, Zn, Zr}$).
    \item \textbf{Pore Accessibility:} Retained systems with Pore Limiting Diameter ($\mathrm{PLD} \ge 2.5\,\text{\AA}$) to permit unhindered diffusion of $\mathrm{H_2O}$ ($2.65\,\text{\AA}$) and $\mathrm{CO_2}$ ($3.30\,\text{\AA}$).
    \item \textbf{Unit Cell Atom Count:} $N_{\mathrm{atoms}} \le 200$ to maintain structural tractability for precision equivariant MLIP relaxation and localized vibrational thermochemistry.
\end{enumerate}

\subsection{Automated Open Metal Site (OMS) and Coordination Vector Algorithm}
For each framework, the coordination sphere around the transition metal nodes was evaluated within a radial cutoff $R = 2.5\,\text{\AA}$. The outward open coordination unit vector $\hat{u}_{\mathrm{open}}$ pointing toward the pore channel was computed as:
\begin{equation}
    \vec{v}_{\mathrm{open}} = -\sum_{i=1}^{\mathrm{CN}} \frac{\vec{r}_{\mathrm{ligand},i} - \vec{r}_{\mathrm{metal}}}{\|\vec{r}_{\mathrm{ligand},i} - \vec{r}_{\mathrm{metal}}\|}, \quad \hat{u}_{\mathrm{open}} = \frac{\vec{v}_{\mathrm{open}}}{\|\vec{v}_{\mathrm{open}}\|}
\end{equation}
Seven reaction intermediates were attached along $\hat{u}_{\mathrm{open}}$:
\begin{itemize}
    \item \textbf{HER:} $*H$ ($d_{\mathrm{M-H}} = 1.55\,\text{\AA}$).
    \item \textbf{OER:} $*OH$ ($1.85\,\text{\AA}$), $*O$ ($1.68\,\text{\AA}$), $*OOH$ ($1.85\,\text{\AA}$, $\angle\mathrm{MOO} = 110^\circ$).
    \item \textbf{CO$_2$RR:} $*COOH$ ($1.95\,\text{\AA}$, C-bound), $*CO$ ($1.85\,\text{\AA}$, C-bound), $*OCHO$ ($1.95\,\text{\AA}$, O-bound).
\end{itemize}

\subsection{Two-Tier MLIP Strategy: CHGNet and MACE-MP-0}
Calculations were orchestrated via the Atomic Simulation Environment (ASE) with CUDA acceleration:
\begin{enumerate}
    \item \textbf{Tier 1 (CHGNet):} Rapid pre-relaxation ($f_{\mathrm{max}} < 0.10\,\text{eV/\AA}$, BFGS) to relieve initial steric contacts and verify framework stability.
    \item \textbf{Tier 2 (MACE-MP-0):} Higher-order equivariant message passing relaxation in double precision ($\mathrm{float64}$, $f_{\mathrm{max}} < 0.03\,\text{eV/\AA}$) to determine ground-state electronic energies $E_{\mathrm{MLIP}}$.
    \item \textbf{Epistemic Uncertainty Filter:} $\sigma_{\mathrm{MLIP}} = |E_{\mathrm{MACE}} - E_{\mathrm{CHGNet}}| / N_{\mathrm{atoms}}$ to identify sterically congested pore pockets.
\end{enumerate}

\subsection{Partial Hessian Vibrational Analysis (PHVA) and CHE Energetics}
Zero-point vibrational energy (ZPE) and vibrational entropy ($S_{\mathrm{vib}}$) at $T = 298.15\,\mathrm{K}$ were evaluated via finite-difference displacements restricted to the adsorbate atoms. Free energies were calculated within the N{\o}rskov Computational Hydrogen Electrode (CHE) framework:
\begin{equation}
    \Delta G = \Delta E_{\mathrm{MLIP}} + \Delta \mathrm{ZPE} - T \Delta S_{\mathrm{vib}} - eU + k_B T \ln(10) \cdot \mathrm{pH}
\end{equation}
Theoretical overpotentials:
\begin{itemize}
    \item $\eta^{\mathrm{HER}} = |\Delta G_{*H}| / e$
    \item $\eta^{\mathrm{OER}} = \max(\Delta G_1, \Delta G_2, \Delta G_3, \Delta G_4)/e - 1.23\,\mathrm{V}$
    \item $\eta^{\mathrm{CO2RR}} = \max(\Delta G_{*COOH}, \Delta G_{*CO} - \Delta G_{*COOH})/e - (-0.11\,\mathrm{V})$
    \item $\Delta G_{\mathrm{sel}} = \Delta G_{*COOH} - \Delta G_{*H}$
\end{itemize}
